\documentclass[10pt,a4paper]{amsart}
\usepackage{amsmath,amssymb,amsthm,graphicx}
% Retained preamble for the Stacks Project chapter snapshot in this directory.
% The chapter's own preamble.tex is a copy of the shared Stacks Project preamble; the block
% below is sliced MECHANICALLY, line for line, from it. Nothing was added, removed or
% redefined; the wrapper class, the xr-hyper/hyperref wiring and the externaldocument list
% are omitted because the assembler supplies the document class.
%
% lines taken from preamble.tex: 3, 5, 9, 11, 12, 156, 157, 158, ... (51 lines total)
%
% package coverage restored from the chapter's own preamble (2026-10-13): xy, xy-2cell

\usepackage[all]{xy}\usepackage{verbatim}
\xyoption{2cell}
\UseAllTwocells

\newenvironment{reference}{\comment}{\endcomment}
\newenvironment{slogan}{\comment}{\endcomment}
\newenvironment{history}{\comment}{\endcomment}
\theoremstyle{plain}
\newtheorem{theorem}[subsection]{Theorem}
\newtheorem{proposition}[subsection]{Proposition}
\newtheorem{lemma}[subsection]{Lemma}
\theoremstyle{definition}
\newtheorem{definition}[subsection]{Definition}
\newtheorem{example}[subsection]{Example}
\newtheorem{exercise}[subsection]{Exercise}
\newtheorem{situation}[subsection]{Situation}
\theoremstyle{remark}
\newtheorem{remark}[subsection]{Remark}
\newtheorem{remarks}[subsection]{Remarks}
\numberwithin{equation}{subsection}
\def\lim{\mathop{\mathrm{lim}}\nolimits}
\def\colim{\mathop{\mathrm{colim}}\nolimits}
\def\Spec{\mathop{\mathrm{Spec}}}
\def\Hom{\mathop{\mathrm{Hom}}\nolimits}
\def\Ext{\mathop{\mathrm{Ext}}\nolimits}
\def\SheafHom{\mathop{\mathcal{H}\!\mathit{om}}\nolimits}
\def\SheafExt{\mathop{\mathcal{E}\!\mathit{xt}}\nolimits}
\def\Sch{\mathit{Sch}}
\def\Mor{\mathop{\mathrm{Mor}}\nolimits}
\def\Ob{\mathop{\mathrm{Ob}}\nolimits}
\def\Sh{\mathop{\mathit{Sh}}\nolimits}
\def\NL{\mathop{N\!L}\nolimits}
\def\CH{\mathop{\mathrm{CH}}\nolimits}
\def\proetale{{pro\text{-}\acute{e}tale}}
\def\etale{{\acute{e}tale}}
\def\QCoh{\mathit{QCoh}}
\def\Ker{\mathop{\mathrm{Ker}}}
\def\Im{\mathop{\mathrm{Im}}}
\def\Coker{\mathop{\mathrm{Coker}}}
\def\Coim{\mathop{\mathrm{Coim}}}
\DeclareMathSymbol{\boxtimes}{\mathbin}{AMSa}{"02}
\def\QCohstack{\mathcal{QC}\!\mathit{oh}}
\def\Cohstack{\mathcal{C}\!\mathit{oh}}
\def\Spacesstack{\mathcal{S}\!\mathit{paces}}
\def\Quotfunctor{\mathrm{Quot}}
\def\Hilbfunctor{\mathrm{Hilb}}
\def\Curvesstack{\mathcal{C}\!\mathit{urves}}
\def\Polarizedstack{\mathcal{P}\!\mathit{olarized}}
\def\Complexesstack{\mathcal{C}\!\mathit{omplexes}}
\def\Pic{\mathop{\mathrm{Pic}}\nolimits}
\def\Picardstack{\mathcal{P}\!\mathit{ic}}
\def\Picardfunctor{\mathrm{Pic}}
\def\Deformationcategory{\mathcal{D}\!\mathit{ef}}
\usepackage{cleveref}
\crefname{originallemma}{original lemma}{original lemmas}
\Crefname{originallemma}{Original lemma}{Original lemmas}
\makeatletter
\newlabel{definition-approximation}{{08HJ}{source}{}{sourcepointer-1}{}}
\newlabel{definition-approximation@cref}{{[originallemma][0][]08HJ}{[1][0][]source}{}{}{}}
\makeatother
\crefname{originallemma}{original lemma}{original lemmas}
\Crefname{originallemma}{Original lemma}{Original lemmas}
\makeatletter
\newlabel{lemma-approximation-affine}{{08HM}{source}{}{sourcepointer-2}{}}
\newlabel{lemma-approximation-affine@cref}{{[originallemma][0][]08HM}{[1][0][]source}{}{}{}}
\makeatother
\crefname{originallemma}{original lemma}{original lemmas}
\Crefname{originallemma}{Original lemma}{Original lemmas}
\makeatletter
\newlabel{lemma-derived-quasi-coherent-small-etale-site}{{071Q}{source}{}{sourcepointer-3}{}}
\newlabel{lemma-derived-quasi-coherent-small-etale-site@cref}{{[originallemma][0][]071Q}{[1][0][]source}{}{}{}}
\makeatother
\crefname{originallemma}{original lemma}{original lemmas}
\Crefname{originallemma}{Original lemma}{Original lemmas}
\makeatletter
\newlabel{lemma-descend-perfect}{{08HG}{source}{}{sourcepointer-4}{}}
\newlabel{lemma-descend-perfect@cref}{{[originallemma][0][]08HG}{[1][0][]source}{}{}{}}
\makeatother
\crefname{originallemma}{original lemma}{original lemmas}
\Crefname{originallemma}{Original lemma}{Original lemmas}
\makeatletter
\newlabel{lemma-descend-pseudo-coherent}{{08HE}{source}{}{sourcepointer-5}{}}
\newlabel{lemma-descend-pseudo-coherent@cref}{{[originallemma][0][]08HE}{[1][0][]source}{}{}{}}
\makeatother
\crefname{originallemma}{original lemma}{original lemmas}
\Crefname{originallemma}{Original lemma}{Original lemmas}
\makeatletter
\newlabel{lemma-glue}{{08HB}{source}{}{sourcepointer-6}{}}
\newlabel{lemma-glue@cref}{{[originallemma][0][]08HB}{[1][0][]source}{}{}{}}
\makeatother
\crefname{originallemma}{original lemma}{original lemmas}
\Crefname{originallemma}{Original lemma}{Original lemmas}
\makeatletter
\newlabel{lemma-open}{{08HL}{source}{}{sourcepointer-7}{}}
\newlabel{lemma-open@cref}{{[originallemma][0][]08HL}{[1][0][]source}{}{}{}}
\makeatother
\crefname{originallemma}{original lemma}{original lemmas}
\Crefname{originallemma}{Original lemma}{Original lemmas}
\makeatletter
\newlabel{perfect-lemma-lift-map}{{08EH}{source}{}{sourcepointer-8}{}}
\newlabel{perfect-lemma-lift-map@cref}{{[originallemma][0][]08EH}{[1][0][]source}{}{}{}}
\makeatother
\crefname{originallemma}{original lemma}{original lemmas}
\Crefname{originallemma}{Original lemma}{Original lemmas}
\makeatletter
\newlabel{perfect-lemma-lift-perfect-complex-plus-shift-support}{{08EK}{source}{}{sourcepointer-9}{}}
\newlabel{perfect-lemma-lift-perfect-complex-plus-shift-support@cref}{{[originallemma][0][]08EK}{[1][0][]source}{}{}{}}
\makeatother
\crefname{originallemma}{original lemma}{original lemmas}
\Crefname{originallemma}{Original lemma}{Original lemmas}
\makeatletter
\newlabel{perfect-section-lift}{{08EC}{source}{}{sourcepointer-10}{}}
\newlabel{perfect-section-lift@cref}{{[originallemma][0][]08EC}{[1][0][]source}{}{}{}}
\makeatother
\crefname{originallemma}{original lemma}{original lemmas}
\Crefname{originallemma}{Original lemma}{Original lemmas}
\makeatletter
\newlabel{remark-match-total-direct-images}{{08GH}{source}{}{sourcepointer-11}{}}
\newlabel{remark-match-total-direct-images@cref}{{[originallemma][0][]08GH}{[1][0][]source}{}{}{}}
\makeatother
\crefname{originallemma}{original lemma}{original lemmas}
\Crefname{originallemma}{Original lemma}{Original lemmas}
\makeatletter
\newlabel{sites-cohomology-lemma-cone-pseudo-coherent}{{08FV}{source}{}{sourcepointer-12}{}}
\newlabel{sites-cohomology-lemma-cone-pseudo-coherent@cref}{{[originallemma][0][]08FV}{[1][0][]source}{}{}{}}
\makeatother
\crefname{originallemma}{original lemma}{original lemmas}
\Crefname{originallemma}{Original lemma}{Original lemmas}
\makeatletter
\newlabel{sites-cohomology-lemma-two-out-of-three-perfect}{{08G9}{source}{}{sourcepointer-13}{}}
\newlabel{sites-cohomology-lemma-two-out-of-three-perfect@cref}{{[originallemma][0][]08G9}{[1][0][]source}{}{}{}}
\makeatother
\title{Stacks, Derived Categories of Spaces: induction-step}
\author{The Stacks Project Authors}
\date{stacks-project:spaces-perfect.tex:local-snapshot-20261012; GNU FDL 1.2 or later}
\begin{document}
\maketitle
\noindent\textit{Editorial note.} Statement span is lines 3684-3691 of sources/stacks-spaces-perfect/pinned-original-source.tex; proof span is lines 3693-3783, the complete author proof (89 body lines).\par\medskip
\noindent\textit{Editorial note.} Navigation stubs: 13 cited-but-not-included labels are supplied as non-load-bearing pointers with their pinned tags where the snapshot tag index has them.\par\medskip
\noindent\textit{Editorial note.} extra\_packages is empty on purpose: the retained preamble already loads xy with options and a second declaration clashes in the assembler wrapper.\par\medskip
\noindent\textit{Editorial note.} Text is retained verbatim, including the author's own typos and placeholders.\par\medskip
\noindent\textit{Editorial note.} Queue index 22 of queue\_spaceperfect.json (pb=3935, stmt 3684-3691, proof 3693-3783).\par\medskip
\noindent\textit{Editorial note.} stacks-spaces-perfect wave sp01 (2026-10-14): fresh adjudication accept per lane-s3/sp01\_judgments; pool per Chief 12:4x unfreeze (reserve-pool workflow, six-check clean); spaces-perfect.tex 126 with-proof rows, 70 b107-occupied labels avoided, 56 free rows judged 32 accept / 24 reject.\par\medskip
\noindent\textit{Editorial note.} Cross-chapter same-label warning (NOT an occupancy conflict, Chief ruling 2026-10-13 evening): label 'lemma-induction-step' also exists in stacks-flat — a different proposition in a different chapter file.\par\medskip
\noindent\textit{Allowed standard prerequisites (editorial).}\begin{itemize}
\item Explicit prerequisite boundary. The cited inputs are prerequisites; the argument in the retained span is not replaced by any cited result.
\end{itemize}\medskip
\ifdefined\hypertarget\hypertarget{sourcepointer-1}{}\fi
\noindent\textit{Source reference.} Original lemma 08HJ: sources/stacks-spaces-perfect/pinned-original-source.tex, source lines 3595--3605. Cited by the retained statement or proof but not part of the collected spans; supplied only as a navigation pointer. The argument does not depend on its content.\par\medskip
\ifdefined\hypertarget\hypertarget{sourcepointer-2}{}\fi
\noindent\textit{Source reference.} Original lemma 08HM: sources/stacks-spaces-perfect/pinned-original-source.tex, source lines 3651--3660. Cited by the retained statement or proof but not part of the collected spans; supplied only as a navigation pointer. The argument does not depend on its content.\par\medskip
\ifdefined\hypertarget\hypertarget{sourcepointer-3}{}\fi
\noindent\textit{Source reference.} Original lemma 071Q: sources/stacks-spaces-perfect/pinned-original-source.tex, source lines 402--407. Cited by the retained statement or proof but not part of the collected spans; supplied only as a navigation pointer. The argument does not depend on its content.\par\medskip
\ifdefined\hypertarget\hypertarget{sourcepointer-4}{}\fi
\noindent\textit{Source reference.} Original lemma 08HG: sources/stacks-spaces-perfect/pinned-original-source.tex, source lines 3288--3293. Cited by the retained statement or proof but not part of the collected spans; supplied only as a navigation pointer. The argument does not depend on its content.\par\medskip
\ifdefined\hypertarget\hypertarget{sourcepointer-5}{}\fi
\noindent\textit{Source reference.} Original lemma 08HE: sources/stacks-spaces-perfect/pinned-original-source.tex, source lines 3133--3139. Cited by the retained statement or proof but not part of the collected spans; supplied only as a navigation pointer. The argument does not depend on its content.\par\medskip
\ifdefined\hypertarget\hypertarget{sourcepointer-6}{}\fi
\noindent\textit{Source reference.} Original lemma 08HB: sources/stacks-spaces-perfect/pinned-original-source.tex, source lines 2206--2213. Cited by the retained statement or proof but not part of the collected spans; supplied only as a navigation pointer. The argument does not depend on its content.\par\medskip
\ifdefined\hypertarget\hypertarget{sourcepointer-7}{}\fi
\noindent\textit{Source reference.} Original lemma 08HL: sources/stacks-spaces-perfect/pinned-original-source.tex, source lines 3628--3636. Cited by the retained statement or proof but not part of the collected spans; supplied only as a navigation pointer. The argument does not depend on its content.\par\medskip
\ifdefined\hypertarget\hypertarget{sourcepointer-8}{}\fi
\noindent\textit{Source reference.} Original lemma 08EH: sources/stacks-spaces-perfect/pinned-original-source.tex, source lines 3783--3783. Cited by the retained statement or proof but not part of the collected spans; supplied only as a navigation pointer. The argument does not depend on its content.\par\medskip
\ifdefined\hypertarget\hypertarget{sourcepointer-9}{}\fi
\noindent\textit{Source reference.} Original lemma 08EK: sources/stacks-spaces-perfect/pinned-original-source.tex, source lines 3783--3783. Cited by the retained statement or proof but not part of the collected spans; supplied only as a navigation pointer. The argument does not depend on its content.\par\medskip
\ifdefined\hypertarget\hypertarget{sourcepointer-10}{}\fi
\noindent\textit{Source reference.} Original lemma 08EC: sources/stacks-spaces-perfect/pinned-original-source.tex, source lines 3783--3783. Cited by the retained statement or proof but not part of the collected spans; supplied only as a navigation pointer. The argument does not depend on its content.\par\medskip
\ifdefined\hypertarget\hypertarget{sourcepointer-11}{}\fi
\noindent\textit{Source reference.} Original lemma 08GH: sources/stacks-spaces-perfect/pinned-original-source.tex, source lines 896--962. Cited by the retained statement or proof but not part of the collected spans; supplied only as a navigation pointer. The argument does not depend on its content.\par\medskip
\ifdefined\hypertarget\hypertarget{sourcepointer-12}{}\fi
\noindent\textit{Source reference.} Original lemma 08FV: sources/stacks-spaces-perfect/pinned-original-source.tex, source lines 3783--3783. Cited by the retained statement or proof but not part of the collected spans; supplied only as a navigation pointer. The argument does not depend on its content.\par\medskip
\ifdefined\hypertarget\hypertarget{sourcepointer-13}{}\fi
\noindent\textit{Source reference.} Original lemma 08G9: sources/stacks-spaces-perfect/pinned-original-source.tex, source lines 3783--3783. Cited by the retained statement or proof but not part of the collected spans; supplied only as a navigation pointer. The argument does not depend on its content.\par\medskip
\begin{lemma}
\label{lemma-induction-step}
Let $S$ be a scheme. Let $(U \subset X, j : V \to X)$ be an
elementary distinguished square of algebraic spaces over $S$.
Assume $U$ quasi-compact, $V$ affine, and $U \times_X V$ quasi-compact.
If approximation by perfect complexes holds on $U$,
then approximation by perfect complexes holds on $X$.
\end{lemma}

\begin{proof}
Let $T \subset |X|$ be a closed subset with $X \setminus T \to X$
quasi-compact. Let $r_U$ be the integer of
Definition \ref{definition-approximation}
adapted to the pair $(U, T \cap |U|)$.
Set $T' = T \setminus |U|$. Endow $T'$ with the induced reduced
subspace structure. Since $|T'|$ is contained in $|X| \setminus |U|$
we see that $j^{-1}(T') \to T'$ is an isomorphism. Moreover,
$V \setminus j^{-1}(T')$ is quasi-compact as it is the fibre product
of $U \times_X V$ with $X \setminus T$ over $X$ and we've assumed
$U \times_X V$ quasi-compact and $X \setminus T \to X$ quasi-compact.
Let $r'$ be the number of affines needed to cover $V \setminus j^{-1}(T')$.
We claim that $r = \max(r_U, r')$ works for the pair $(X, T)$.

\medskip\noindent
To see this choose a triple $(T, E, m)$ such that $E$ is
$(m - r)$-pseudo-coherent and $H^i(E)$ is supported on $T$ for
$i \geq m - r$. Let $t$ be the largest integer such that
$H^t(E)|_U$ is nonzero. (Such an integer exists as $U$ is quasi-compact
and $E|_U$ is $(m - r)$-pseudo-coherent.)
We will prove that $E$ can be approximated by induction on $t$.

\medskip\noindent
Base case: $t \leq m - r'$. This means that $H^i(E)$ is supported
on $T'$ for $i \geq m - r'$. Hence
Lemma \ref{lemma-approximation-affine}
guarantees the existence of an approximation
$P \to E|_V$ of $(T', E|_V, m)$ on $V$.
Applying Lemma \ref{lemma-open} we see that
$(T', E, m)$ can be approximated. Such an approximation
is also an approximation of $(T, E, m)$.

\medskip\noindent
Induction step. Choose an approximation $P \to E|_U$ of
$(T \cap |U|, E|_U, m)$. This in particular gives a surjection
$H^t(P) \to H^t(E|_U)$.
In the rest of the proof we will use the equivalence of
Lemma \ref{lemma-derived-quasi-coherent-small-etale-site}
(and the compatibilities of Remark \ref{remark-match-total-direct-images})
for the representable algebraic spaces $V$ and $U \times_X V$.
We will also use the fact that $(m - r)$-pseudo-coherence,
resp.\ perfectness on the Zariski site and \'etale site agree, see
Lemmas \ref{lemma-descend-pseudo-coherent} and
\ref{lemma-descend-perfect}.
Thus we can use the results of
Derived Categories of Schemes, Section \ref{perfect-section-lift}
for the open immersion $U \times_X V \subset V$. In this way
Derived Categories of Schemes,
Lemma \ref{perfect-lemma-lift-perfect-complex-plus-shift-support}
implies there exists a perfect object $Q$ in $D(\mathcal{O}_V)$
supported on $j^{-1}(T)$ and an isomorphism
$Q|_{U \times_X V} \to (P \oplus P[1])|_{U \times_X V}$. By
Derived Categories of Schemes, Lemma \ref{perfect-lemma-lift-map}
we can replace $Q$ by $Q \otimes^\mathbf{L} I$
and assume that the map
$$
Q|_{U \times_X V} \longrightarrow
(P \oplus P[1])|_{U \times_X V} \longrightarrow
P|_{U \times_X V} \longrightarrow E|_{U \times_X V}
$$
lifts to $Q \to E|_V$. By Lemma \ref{lemma-glue}
we find an morphism $a : R \to E$ of $D(\mathcal{O}_X)$
such that $a|_U$ is isomorphic to $P \oplus P[1] \to E|_U$
and $a|_V$ isomorphic to $Q \to E|_V$.
Thus $R$ is perfect and supported on $T$
and the map $H^t(R) \to H^t(E)$ is surjective on restriction to $U$.
Choose a distinguished triangle
$$
R \to E \to E' \to R[1]
$$
Then $E'$ is $(m - r)$-pseudo-coherent
(Cohomology on Sites, Lemma \ref{sites-cohomology-lemma-cone-pseudo-coherent}),
$H^i(E')|_U = 0$ for $i \geq t$, and
$H^i(E')$ is supported on $T$ for $i \geq m - r$.
By induction we find an approximation $R' \to E'$
of $(T, E', m)$. Fit the composition $R' \to E' \to R[1]$
into a distinguished triangle $R \to R'' \to R' \to R[1]$
and extend the morphisms $R' \to E'$ and $R[1] \to R[1]$ into
a morphism of distinguished triangles
$$
\xymatrix{
R \ar[r] \ar[d] & R'' \ar[d] \ar[r] & R' \ar[d] \ar[r] & R[1] \ar[d] \\
R \ar[r] &  E \ar[r] & E' \ar[r] & R[1]
}
$$
using TR3. Then $R''$ is a perfect complex
(Cohomology on Sites, Lemma
\ref{sites-cohomology-lemma-two-out-of-three-perfect})
supported on $T$. An easy diagram chase shows that $R'' \to E$ is the desired
approximation.
\end{proof}

\end{document}
